\documentclass[10pt,a4paper]{amsart}
\usepackage{amsmath,amssymb,amsthm,graphicx}
\usepackage{verbatim}
\newenvironment{reference}{\comment}{\endcomment}
%\newenvironment{reference}{}{}
\newenvironment{slogan}{\comment}{\endcomment}
\newenvironment{history}{\comment}{\endcomment}

% For commutative diagrams we use Xy-pic
\usepackage[all]{xy}

% We use 2cell for 2-commutative diagrams.
\xyoption{2cell}
% Theorem environments.
%
\theoremstyle{plain}
\newtheorem{theorem}[subsection]{Theorem}
\newtheorem{proposition}[subsection]{Proposition}
\newtheorem{lemma}[subsection]{Lemma}

\theoremstyle{definition}
\newtheorem{definition}[subsection]{Definition}
\newtheorem{example}[subsection]{Example}
\newtheorem{exercise}[subsection]{Exercise}
\newtheorem{situation}[subsection]{Situation}

\theoremstyle{remark}
\newtheorem{remark}[subsection]{Remark}
\newtheorem{remarks}[subsection]{Remarks}

\numberwithin{equation}{subsection}

% Macros
%
\def\lim{\mathop{\mathrm{lim}}\nolimits}
\def\colim{\mathop{\mathrm{colim}}\nolimits}
\def\Spec{\mathop{\mathrm{Spec}}}
\def\Hom{\mathop{\mathrm{Hom}}\nolimits}
\def\Ext{\mathop{\mathrm{Ext}}\nolimits}
\def\SheafHom{\mathop{\mathcal{H}\!\mathit{om}}\nolimits}
\def\SheafExt{\mathop{\mathcal{E}\!\mathit{xt}}\nolimits}
\def\Sch{\mathit{Sch}}
\def\Mor{\mathop{\mathrm{Mor}}\nolimits}
\def\Ob{\mathop{\mathrm{Ob}}\nolimits}
\def\Sh{\mathop{\mathit{Sh}}\nolimits}
\def\NL{\mathop{N\!L}\nolimits}
\def\CH{\mathop{\mathrm{CH}}\nolimits}
\def\proetale{{pro\text{-}\acute{e}tale}}
\def\etale{{\acute{e}tale}}
\def\QCoh{\mathit{QCoh}}
\def\Ker{\mathop{\mathrm{Ker}}}
\def\Im{\mathop{\mathrm{Im}}}
\def\Coker{\mathop{\mathrm{Coker}}}
\def\Coim{\mathop{\mathrm{Coim}}}

% Boxtimes
%
\DeclareMathSymbol{\boxtimes}{\mathbin}{AMSa}{"02}

%
% Macros for moduli stacks/spaces
%
\def\QCohstack{\mathcal{QC}\!\mathit{oh}}
\def\Cohstack{\mathcal{C}\!\mathit{oh}}
\def\Spacesstack{\mathcal{S}\!\mathit{paces}}
\def\Quotfunctor{\mathrm{Quot}}
\def\Hilbfunctor{\mathrm{Hilb}}
\def\Curvesstack{\mathcal{C}\!\mathit{urves}}
\def\Polarizedstack{\mathcal{P}\!\mathit{olarized}}
\def\Complexesstack{\mathcal{C}\!\mathit{omplexes}}
% \Pic is the operator that assigns to X its picard group, usage \Pic(X)
% \Picardstack_{X/B} denotes the Picard stack of X over B
% \Picardfunctor_{X/B} denotes the Picard functor of X over B
\def\Pic{\mathop{\mathrm{Pic}}\nolimits}
\def\Picardstack{\mathcal{P}\!\mathit{ic}}
\def\Picardfunctor{\mathrm{Pic}}
\def\Deformationcategory{\mathcal{D}\!\mathit{ef}}
\usepackage{cleveref}
\crefname{originallemma}{original lemma}{original lemmas}
\Crefname{originallemma}{Original lemma}{Original lemmas}
\makeatletter
\newlabel{lemma-functor-on-fp-modules}{{0GNJ}{source}{}{sourcepointer-1}{}}
\newlabel{lemma-functor-on-fp-modules@cref}{{[originallemma][0][]0GNJ}{[1][0][]source}{}{}{}}
\makeatother
\crefname{originalreference}{original reference}{original references}
\Crefname{originalreference}{Original reference}{Original references}
\makeatletter
\newlabel{algebra-definition-coherent}{{05CV}{source}{}{sourcepointer-2}{}}
\newlabel{algebra-definition-coherent@cref}{{[originalreference][0][]05CV}{[1][0][]source}{}{}{}}
\makeatother
\crefname{originalreference}{original reference}{original references}
\Crefname{originalreference}{Original reference}{Original references}
\makeatletter
\newlabel{algebra-lemma-coherent-ring}{{05CX}{source}{}{sourcepointer-3}{}}
\newlabel{algebra-lemma-coherent-ring@cref}{{[originalreference][0][]05CX}{[1][0][]source}{}{}{}}
\makeatother
\crefname{originalreference}{original reference}{original references}
\Crefname{originalreference}{Original reference}{Original references}
\makeatletter
\newlabel{algebra-lemma-coherent}{{05CW}{source}{}{sourcepointer-4}{}}
\newlabel{algebra-lemma-coherent@cref}{{[originalreference][0][]05CW}{[1][0][]source}{}{}{}}
\makeatother
\crefname{originalreference}{original reference}{original references}
\Crefname{originalreference}{Original reference}{Original references}
\makeatletter
\newlabel{categories-definition-exact}{{0034}{source}{}{sourcepointer-5}{}}
\newlabel{categories-definition-exact@cref}{{[originalreference][0][]0034}{[1][0][]source}{}{}{}}
\makeatother
\crefname{originalreference}{original reference}{original references}
\Crefname{originalreference}{Original reference}{Original references}
\makeatletter
\newlabel{homology-lemma-additive-functor}{{0DLP}{source}{}{sourcepointer-6}{}}
\newlabel{homology-lemma-additive-functor@cref}{{[originalreference][0][]0DLP}{[1][0][]source}{}{}{}}
\makeatother
\crefname{originallemma}{original lemma}{original lemmas}
\Crefname{originallemma}{Original lemma}{Original lemmas}
\makeatletter
\newlabel{lemma-functor-on-fp-modules-additive}{{0GNK}{source}{}{sourcepointer-7}{}}
\newlabel{lemma-functor-on-fp-modules-additive@cref}{{[originallemma][0][]0GNK}{[1][0][]source}{}{}{}}
\makeatother
\crefname{originalreference}{original reference}{original references}
\Crefname{originalreference}{Original reference}{Original references}
\makeatletter
\newlabel{categories-lemma-characterize-left-exact}{{0035}{source}{}{sourcepointer-8}{}}
\newlabel{categories-lemma-characterize-left-exact@cref}{{[originalreference][0][]0035}{[1][0][]source}{}{}{}}
\makeatother
\crefname{originalreference}{original reference}{original references}
\Crefname{originalreference}{Original reference}{Original references}
\makeatletter
\newlabel{algebra-lemma-module-colimit-fp}{{00HA}{source}{}{sourcepointer-9}{}}
\newlabel{algebra-lemma-module-colimit-fp@cref}{{[originalreference][0][]00HA}{[1][0][]source}{}{}{}}
\makeatother
\title{Left exactness of the filtered-colimit extension over a coherent ring}
\author{The Stacks Project Authors}
\date{stacks-project:functors.tex:local-snapshot-20261012; GNU FDL 1.2 or later}
\begin{document}
\maketitle
\noindent\textit{Editorial note.} All delivered mathematical text is the author's own, copied verbatim from the pinned Stacks Project chapter file functors.tex; no statement or proof step was rewritten, reordered or removed.\par\medskip
\noindent\textit{Editorial note.} Span recomputation by label pairing (LANE\_BRIEF 10.2): in this chapter the author's proof environment follows the statement environment as a standalone block; the spec covers the complete environment pair (statement span, proof span) as recorded in the pool candidate.\par\medskip
\noindent\textit{Editorial note.} Occupation scan: combined claim\_key + statement\_anchor + generic-line scan over all registration files (active and withdrawn) found 0 hits for functors.tex (fu-occupied.json, 2026-10-14); consistent with the duty manager's six-check dispatch.\par\medskip
\ifdefined\hypertarget\hypertarget{sourcepointer-1}{}\fi
\noindent\textit{Source reference.} Original lemma 0GNJ: sources/stacks-chow/tags.txt, source lines 20794--20794. Cross-chapter standard prerequisite cited by the author inside the counted span; it is used as a pointer and its own proof is not part of this item.\par\medskip
\ifdefined\hypertarget\hypertarget{sourcepointer-2}{}\fi
\noindent\textit{Source reference.} Original reference 05CV: sources/stacks-chow/tags.txt, source lines 6933--6933. Cross-chapter standard prerequisite cited by the author inside the counted span; it is used as a pointer and its own proof is not part of this item.\par\medskip
\ifdefined\hypertarget\hypertarget{sourcepointer-3}{}\fi
\noindent\textit{Source reference.} Original reference 05CX: sources/stacks-chow/tags.txt, source lines 6935--6935. Cross-chapter standard prerequisite cited by the author inside the counted span; it is used as a pointer and its own proof is not part of this item.\par\medskip
\ifdefined\hypertarget\hypertarget{sourcepointer-4}{}\fi
\noindent\textit{Source reference.} Original reference 05CW: sources/stacks-chow/tags.txt, source lines 6934--6934. Cross-chapter standard prerequisite cited by the author inside the counted span; it is used as a pointer and its own proof is not part of this item.\par\medskip
\ifdefined\hypertarget\hypertarget{sourcepointer-5}{}\fi
\noindent\textit{Source reference.} Original reference 0034: sources/stacks-chow/tags.txt, source lines 121--121. Cross-chapter standard prerequisite cited by the author inside the counted span; it is used as a pointer and its own proof is not part of this item.\par\medskip
\ifdefined\hypertarget\hypertarget{sourcepointer-6}{}\fi
\noindent\textit{Source reference.} Original reference 0DLP: sources/stacks-chow/tags.txt, source lines 17054--17054. Cross-chapter standard prerequisite cited by the author inside the counted span; it is used as a pointer and its own proof is not part of this item.\par\medskip
\ifdefined\hypertarget\hypertarget{sourcepointer-7}{}\fi
\noindent\textit{Source reference.} Original lemma 0GNK: sources/stacks-chow/tags.txt, source lines 20795--20795. Cross-chapter standard prerequisite cited by the author inside the counted span; it is used as a pointer and its own proof is not part of this item.\par\medskip
\ifdefined\hypertarget\hypertarget{sourcepointer-8}{}\fi
\noindent\textit{Source reference.} Original reference 0035: sources/stacks-chow/tags.txt, source lines 122--122. Cross-chapter standard prerequisite cited by the author inside the counted span; it is used as a pointer and its own proof is not part of this item.\par\medskip
\ifdefined\hypertarget\hypertarget{sourcepointer-9}{}\fi
\noindent\textit{Source reference.} Original reference 00HA: sources/stacks-chow/tags.txt, source lines 636--636. Cross-chapter standard prerequisite cited by the author inside the counted span; it is used as a pointer and its own proof is not part of this item.\par\medskip
\begin{lemma}
\label{lemma-functor-on-fp-modules-left-exact}
Let $A$, $\mathcal{B}$, $F$ be as in Lemma \ref{lemma-functor-on-fp-modules}.
Assume $A$ is a coherent ring (Algebra, Definition
\ref{algebra-definition-coherent}), $\mathcal{B}$ is additive, has kernels,
filtered colimits commute with taking kernels, and $F$ is left exact. Then
$F'$ is additive, left exact, and commutes with arbitrary direct sums.
\end{lemma}

\begin{proof}
Since $A$ is coherent, the category $\text{Mod}^{fp}_A$ is abelian
with same kernels and cokernels as in $\text{Mod}_A$, see
Algebra, Lemmas \ref{algebra-lemma-coherent-ring} and
\ref{algebra-lemma-coherent}. Hence all finite limits exist in
$\text{Mod}^{fp}_A$ and
Categories, Definition \ref{categories-definition-exact} applies.
Since $F$ is left exact, $F$ commutes with products of pairs, which are
represented by direct sums. Hence $F$ is additive by
Homology, Lemma \ref{homology-lemma-additive-functor}.
Hence $F'$ is additive and commutes with direct sums by
Lemma \ref{lemma-functor-on-fp-modules-additive}.
We urge the reader to prove that $F'$ is left exact themselves
instead of reading the proof below.

\medskip\noindent
To show that $F'$ is left exact, it suffices to show that $F'$ commutes
with equalizers, see
Categories, Lemma \ref{categories-lemma-characterize-left-exact}.
Now, if $a, b : L \to M$ are maps of $A$-modules, then the
equalizer of $a$ and $b$ is the kernel of $a - b : L \to M$.
Thus let $0 \to K \to L \to M$ be an exact sequence
of $A$-modules. We have to show that in
$$
0 \to F'(K) \to F'(L) \to F'(M)
$$
the arrow $F'(K) \to F'(L)$ is a kernel for $F'(L) \to F'(M)$ in $\mathcal{B}$
(if $\mathcal{B}$ were abelian we would say that the displayed sequence
is exact).
Write $M = \colim_{i \in I} M_i$ as a filtered colimit of finitely
presented $A$-modules, see
Algebra, Lemma \ref{algebra-lemma-module-colimit-fp}.
Let $L_i = L \times_M M_i$. 
We obtain a system of exact sequences $0 \to K \to L_i \to M_i$
over $I$. Since filtered colimits commute with taking kernels
in $\mathcal{B}$ by assumption, 
it suffices to show that $F'(K) \to F'(L_i)$ is a kernel of
$F'(L_i) \to F(M_i)$ in $\mathcal{B}$ for all $i \in I$. In other words, we may
assume $M$ is finitely presented. Write $L = \colim_{i \in I} L_i$
as a filtered colimit of finitely presented $A$-modules.
Let $K_i = K \times_L L_i$. We obtain a system of short exact sequences
$0 \to K_i \to L_i \to M$ over $I$.
Repeating the argument already given, we reduce to showing
$F'(K_i) \to F(L_i)$ is a kernel of
$F(L_i) \to F(M)$ in $\mathcal{B}$ for all $i \in I$.
In other words, we may assume both
$L$ and $M$ are finitely presented $A$-modules.
Since $A$ is coherent, the $A$-module $K = \Ker(L \to M)$ is
of finite presentation as the category of finitely presented
$A$-modules is abelian (see references given above).
In other words, all three modules $K$, $L$, and $M$
are finitely presented $A$-modules. This final case follows
from the assumption that $F$ is left exact.
\end{proof}

\end{document}
